\section{The Metaphysics of Imperfection}

\emph{One's circle is never closed, one's past is not the future, one's work is never finished ... }

\begin{quotex}
... the primitive craftsman leaves in his work something unfinished, and that the primitive mother dislikes to hear the beauty of her child unduly praised; it is ``tempting Providence'', and may lead to disaster. That seems like nonsense to us. And yet there survives in our vernacular the explanation for the principle involved: the craftsman leaves something undone in his work for the same reason that the word ``to be finished'' may mean either to be perfected or to die. Perfection is death: when a thing has been altogether fulfilled, when all has been done that was to be done, potentiality altogether reduced to act, that is the end: those whom the gods love die young. This is not what the workman desired for his work, nor the mother for her child. It can very well be that the workman or the peasant mother is no longer conscious of the meaning of a precaution that may have become a mere superstition; but assuredly we, who call ourselves anthropologists, should have been able to understand what was the idea which alone could have given rise to such a superstition, and ought to have asked ourselves whether or not the peasant by his actual observance of the precaution is not defending himself from a dangerous suggestion to which we, who have made of our existence a more tightly closed system, may be immune.
\flright{\textsc{Ananda Commaraswamy}, \emph{Primitive Mentality}\url{http://www.gornahoor.net/library/PrimitiveMentality.htm}}
\end{quotex}


This, too, is the practice of the metaphysician. No matter how deep is his understanding or broad his knowledge, there is always something undone: something else to learn, another connection to see, another experience to integrate. The idea that someone can believe in a closed system that encompasses all knowledge, or that the repetition of old practices is totally sufficient, or that the recovery of old beliefs and rites will be the guide for the future are all temptations that too often ``lead to disaster''.

Guenon points out that death comes when all one's possibilities have been exhausted. Even if that death is not physical, it will be spiritual. For as soon as we believe we are spiritually perfect, then we become spiritually dead, no longer able to grow, thrive or pass on an inheritance. The metaphysician is convinced that the Absolute is Infinite, and thus is in no danger of running out of possibilities.


\flrightit{Posted on 2010-07-05 by Cologero}

\begin{center}* * *\end{center}

\begin{footnotesize}
\begin{sffamily}
\texttt{Arturo Vasqiez on 2010-07-05 at 23:11 said:}

In Mexico, if you praise an infant for its beauty, you are supposed to to it, lest you give it the evil eye, and it falls ill as a result. This is very common in many cultures, to the point that children were intentionally scared at times to make them less attractive. There is a certain fatalism at the heart of the cosmos which modern man tends to deny.

\hfill

\texttt{Cologero on 2010-07-05 at 23:49 said:}

\begin{quote}\footnotesize

There is a certain fatalism at the heart of the cosmos which modern man tends to deny.
\end{quote}


That is probably why modern man praises the imperfect and the unworthy ... he considers that fatalism to be an injustice.


\hfill

\texttt{Arturo Vasquez on 2010-07-06 at 07:59 said:}

I would rather attribute the praise of the weak and the victim to Christianity, as would Rene Girard. In Christian societies, the victim is always the moral superior of the victor. However, in our own society, this has gottn to a morally grotesque point, where everyone would play the victim out of a sheer sense of social guile. The worst case of this is of course the whole idea of ``white man victimhood'', but that is another conversation entirely.

\hfill

\texttt{Cologero on 2010-07-06 at 08:36 said:}

I find it curious that you would misread ``the imperfect and the unworthy'' as ``the weak and the victim''. First of all, this blog has often written about the code of chivalry and its obligation to protect the weak (widows, orphans) and the victim of injustice. Apparently I need to emphasize ``protect'' and not ``praise''; the weak and the victim may deserve protection, but why praise? Quite the opposite: it is the knight --- the protector and the bringer of justice --- whom we praise.

Now to get back to the point I actually did make. The modern world praises the morally imperfect and the spiritually weak. It regards the morally praiseworthy and the spiritually strong as an affront and
\emph{ipso facto} an implicit condemnation. Rather than facing up to and correcting one's faults, the modern world, instead, would praise the faulty and condemn the faultless. It regards any inequality as an injustice --- the result of a fatalism (a heathen notion) that must be rejected, rather than as the evidence of Providence. The Church (at least prior to V II) never condemned inequality nor did it ever advocate the right to do wrong.

\hfill

\texttt{Will on 2010-07-06 at 11:54 said:}

Well said, Cologero. The line about the difference between protecting vs. praising the weak highlights one of the inversions of the modern world which you wrote of earlier. I think that, in his own way, this is part of what Nietzsche was getting at in his critique of Christianity. For all his praise of might and ferocity, he also speaks of the proper disposition of the strong and noble man as being one of courtesy and kindness towards the weak and unfortunate.

The analogy of society and the soul from Plato's Republic is also relevant here. The distinction between the strong and weak is somewhat relative. For example, Stephen Hawking is strong in terms of his mathematical ability, but weak in athletic ability. Everyone has their strengths and weaknesses, their talents and shortcomings. So the notion that the strong and noble should be kind to, and protective of, the weaker not only refers to interaction between individuals, but interactions within individuals -- the strong parts of the soul showing kindness and understanding to the weak parts, and thus sharing their strength and working to remedy the imbalance.

I think this is the societal ethic that the Christian teachings point to, but in modern times it has become what Nietzsche called the `society of nurses' -- the weak caring for the weak, and banishing all strength. That doesn't work, neither in society nor in the psyche.

Also, `kindness' here does not necessarily mean `niceness' and tolerant acceptance. Buddhists distinguish between real compassion -- the willingness to do what is necessary based on clear vision and love -- and `idiot compassion' -- the mistaken belief that compassion is simply being nice to everyone all the time.


\hfill

\end{sffamily}
\end{footnotesize}

